\uuid{pzt9}
\titre{Ensemble}
\auteur{bodin ; hanani ; mzari}
\organisation{exo7}
\licence{CC BY-NC-SA 4.0}
\datecreate{2020-01-20}
\niveau{L1}
\module{Algèbre}
\chapitre{Logique, ensemble, raisonnement}
\sousChapitre{Ensemble}
\difficulte{2}
\theme{Ensembles, applications, ensemble}
\exo7qcmid{lille-1.ensembles.015}

\contenu{
\begin{question}
\qclassification{100.02,101.01,102.02,102.02}
\qtags{motcle=ensemble}
Pour tout $n\in \Nn ^*$, on pose $E_n=\{1,2,\dots ,n\}$. On note $\mathcal{P}(E_n)$ l'ensemble des parties de $E_n$. Quelles sont les bonnes réponses ?
\begin{answers}
    \bad{$\mathcal{P}(E_2)=\{\{1\},\{2\}\}$}
    \good{$\mathcal{P}(E_2)=\{\varnothing ,\{1\},\{2\},E_2\}$}
    \bad{$\mathrm{Card}(\mathcal{P}(E_n))=n$}
    \good{$\mathrm{Card}(\mathcal{P}(E_n))=2^n$}
\end{answers}
\begin{explanations}
Le nombre de parties à $k$ éléments de $E_n$ est $\mathrm{C}^k_n$ et le nombre de toutes les parties de $E_n$ est $\displaystyle \sum _{k=0}^n\mathrm{C}^k_n=(1+1)^n=2^n$.
\end{explanations}
\end{question}
}
